(1) 在固定T,V的理想气体容器中发生A+B可逆生成C，初始各物质的量分别a0,b0,c0，反应进度xi。用理想气体Helmholtz自由能及化学势写出平衡反应商Q和无量纲平衡常数K(T)的关系。
(2) 证明平衡自由能对xi严格凸，从而在可行化学组成区间内若存在平衡根，则根唯一；给出初始a0=b0=n,c0=0时关于xi的代数方程。
(3) 若反应标准摩尔焓DeltaH°近似与温度无关，推导van’t Hoff公式d ln K/dT=DeltaH°/(R T²)，分析放热反应的温度趋势。
(4) 独立从反应亲和力和自由能变化说明平衡附近演化方向，核验体积对反应平衡的影响及气体非理想性边界。
